\section{The common field, its integral lattice and its translates}\label{host:sec}

Fix a primitive Dirichlet character \(\psi\) of conductor \(M\), of order
\(o\). The arithmetic field itself and the finite sets used in its
translates admit explicit constructions. Both matter in applying the
prime bounds: changing the field after the pair is chosen could change
the list of twists that was tested.

\begin{proposition}[Cyclotomic realization]\label{host:art}
Put \(N_e=\operatorname{lcm}(2^{e+2},M)\), \(U_e=(\mathbf Z/N_e)^\times\),
\(C_e=(\mathbf Z/2^{e+2})^\times/\{\pm1\}\), \(B=\operatorname{im}\psi\).
Let \(\eta_e:U_e\to C_e\times B\) send \(u\) to \(([u],\psi(u))\),
and set \(H_e=\operatorname{im}\eta_e,\ A_0=\ker(H_e\to B)\).
Then \(A_0\) is cyclic, \([H_e:A_0]=o\), and \(|H_e|=c\,2^e\) for all
sufficiently large \(e\), where \(1\le c\le o\) is independent of \(e\).

Take \(e\ge2\) and a pair \(q,d\) from Theorem~\ref{pr:pair}, with
\(Y^{\varepsilon/5}\ge4M\). There is an abelian CM field \(E\) with
Galois group \(\{\pm1\}^2\times H_e\). The characters trivial on
\(\{(1,1,h):h\in A_0\}\), pulled back to integers prime to \(N_e qd\),
are precisely
\[
      \psi^a J_q^{\,i}J_d^{\,j},
 \qquad 0\le a<o,\quad i,j\in\{0,1\},\qquad J_p(u)=(u/p).
\]
The retained characters are the \(o\) primitive inducers of the
\(\psi^a\). Each remaining one is among the twists tested on the same pair.
\end{proposition}
\begin{proof}
Reduction of unit groups modulo a divisor is onto, by the Chinese
remainder theorem prime by prime. It follows that both projections of
\(H_e\) are onto. The kernel of the second injects into \(C_e\), cyclic
of order \(2^e\): for \(e\ge1\) the class of \(5\) generates, since
\(v_2(5^{2^j}-1)=j+2\) by successive squaring and all its powers are
\(1\bmod4\). For \(e=0\) the quotient is trivial. The kernel of the first
projection identifies with \(B_e=\psi(\{u\in U_e:u\equiv\pm1\bmod
2^{e+2}\})\). Reduction from level \(e+1\) shows \(B_{e+1}\subseteq B_e\).
This descending chain in the fixed finite group \(B\) stabilizes. Take
\(c\) equal to the eventual order and use the exact sequence with quotient
\(C_e\).

Define \(w:\{\pm1\}\to H_e\) by \(w(h)=\eta_e(h)\).
By the Chinese remainder theorem the homomorphism
\[
    u\longmapsto
    \bigl(J_q(u),J_d(u),\,\eta_e(u)w(J_q(u))\bigr)
       \quad\text{on }(\mathbf Z/N_eqd)^\times
\]
is onto. Indeed the signs can be prescribed at \(q,d\), and one then
chooses the \(\eta_e\)-image to attain the specified third coordinate.
Take \(E\) to be the subfield of \(\mathbf Q(\zeta_{N_e qd})\) fixed
by the kernel. Its kernel condition is exactly
\[
 u\equiv\pm1\bmod 2^{e+2},\qquad \psi(u)=J_q(u)=J_d(u)=1.
\]
Thus \(E\) is the compositum of the maximal real \(2^{e+2}\)-cyclotomic
field, the field cut out by \(\psi\), and the quadratic fields of
discriminants \(-q,-d\). In fact \(q,d\equiv3\bmod4\).
Complex conjugation has coordinates \((-1,-1,1)\), since \(w(-1)^2=1\).
Its fixed field is totally real by abelianness, and \(E\) contains the
imaginary quadratic field \(K=\mathbf Q(\sqrt{-d})\).
The integer \(N_e\) is also the least divisor of \(N_e\) through which
\(\psi\) factors and which is divisible by \(2^{e+2}\): a modulus for
the primitive \(\psi\) must be divisible by \(M\).

A character of \(H_e\) killing \(A_0\) is a character of \(B\). Since
\(B\subset\mathbf C^\times\) is cyclic of order \(o\), it raises values
of \(\psi\) to a power. Its value on the \(w(J_q)\) term is a sign
character of \(J_q\), absorbed into \(i\). Conversely choose the character
on the first sign to undo this additional sign. This proves equality of
the two character lists. They consist of \(4o\) distinct characters:
one can vary \(J_q\) and \(J_d\) independently while \(u\bmod N_e\) is
fixed, and conversely can set both signs to one. Passing to primitive
inducers at a fixed common modulus preserves distinctness. In particular
there are exactly three nonempty binary choices tested per power of \(\psi\).
\end{proof}

\begin{lemma}[The half-different]\label{host:different}
There is an integral ideal \(\mathfrak b\) of \(E\), invariant under
\(\operatorname{Gal}(E/K)\), such that
\(\mathfrak D_{E/K}=\mathfrak b\,\overline{\mathfrak b}\) and
\(\mathfrak b+\overline{\mathfrak b}=\mathcal O_E\).
The \(\mathcal O_K\)-lattice \(I=\mathfrak b^{-1}\) is self-dual for
\((u,v)\mapsto\operatorname{Tr}_{E/K}(u\bar v)\).
\end{lemma}
\begin{proof}
All finite ramification in \(E/K\) lies above primes dividing \(N_eq\).
Indeed the other three factors in the compositum with \(K\) are
unramified outside \(N_eq\), and unramified local extensions stay
unramified after base change. Each such rational prime splits in \(K\).
For odd primes dividing \(M\), and for \(q\), this follows from
\(-d\equiv1\bmod p\). For \(p=2\) use the odd discriminant
\(-d\equiv1\bmod8\). The different is invariant under every
\(\mathbf Q\)-automorphism of \(E\), by the defining trace dual.
For each rational prime in its support choose one of the two places of
\(K\) and place in \(\mathfrak b\) every factor of the different above
that place with its full multiplicity. Conjugation supplies exactly
the disjoint remaining factors; automorphisms fixing \(K\) preserve
the choice. For any fractional ideal \(J\), the integrality of
\(\operatorname{Tr}_{E/K}(u\bar J)\) is equivalent to
\(u\bar J\subseteq\mathfrak D_{E/K}^{-1}\), since \(\bar J\) is closed
under multiplication by \(\mathcal O_E\). Hence
\(I^\#=(\bar I)^{-1}\mathfrak D_{E/K}^{-1}=\mathfrak b^{-1}\).
\end{proof}

\begin{lemma}[The rank fits the cohomology bound]\label{host:rank}
Suppose \(x^{2+8\varepsilon}\le n=|H_e|\le2x^{2+8\varepsilon}\) and
\(0<\varepsilon<1/52\). Put \(f=2n,\ L=x^\varepsilon\) and choose \(m,k\)
to be the least even integers at least \(x^{1/2-\varepsilon/16}\) and
\(x^{9/2+10\varepsilon}\), respectively. Define
\[
 \begin{split}
 r&=k+m-1,\qquad P=\lceil x^{377}\rceil,\\
 D&=8\left\lceil \frac{10r}{8}
       \left(f+\frac{f^2(k-1)^3L^2}{m}\right)\right\rceil,\qquad
 t=8fr+4+PD-r .
 \end{split}
\]
Then \(t\) is a positive odd integer and \(t\le x^{400}\) for all
sufficiently large \(x\), uniformly in \(0<\varepsilon<1/52\).
\end{lemma}
\begin{proof}
Here \(r\) is odd, \(f\ge1\), and \(D\) is divisible by eight, proving
positivity and parity. With \(u=x^{1/100}\ge2\), the stated ranges give
\[
 f\le4u^{216},\quad k\le2u^{470},\quad r\le4u^{470},\quad
 m\ge u^{49},\quad L^2\le u^4.
\]
The rounding bound \(2\lceil v/2\rceil\le2v\) holds for \(v\ge2\).
The exponent of \(f^2(k-1)^3L^2/m\) in \(u\) is at most
\(432+1410+4-49=1797\), with coefficient \(128\). Consequently
\(D\le5288u^{2267}\), while \(P\le2u^{37700}\). Thus
\[
 t\le 10708u^{39967}\le u^{40000}
\]
because \(10708\le2^{33}\le u^{33}\). This strict margin is needed
when the cohomology window is read at the rescaled value in
Section~\ref{sup:section}.
\end{proof}

\begin{proposition}[A dense set with bounded orbit intersections]\label{host:behrend}
For \(n\ge1,\ L\ge2\) integers there is \(S\subset[0,n/(4L))\cap\mathbf Z\)
with
\[
 |S|\ge (n/L)\exp\{-C\sqrt{\log n\,\log L}\}
\]
for an absolute \(C\). For distinct \(s_1,s_2,s_3\in S\) and distinct
nonzero \(i,j\in[-L,L]\cap\mathbf Z\) one has
\(j(s_2-s_1)\ne i(s_3-s_1)\). In \(\mathbf Z/n\), every translate of
\(S\) meets \(L\) consecutive terms of an orbit of step \(h\) in at most
two terms if \(\operatorname{ord}(h)>4L\). For
\(\operatorname{ord}(h)\le4L\) it meets the whole orbit in at most one
distinct point.
\end{proposition}
\begin{proof}
Set \(v=\log(4L)\). If \(\log n\le64v\), the singleton \(\{0\}\)
suffices on increasing the absolute constant, since \(v\le3\log L\).
Otherwise put \(S_0=\log(n/(4L))\),
\(d_0=\lfloor\sqrt{S_0/v}\rfloor,\ b=\lfloor e^{S_0/d_0}/(4L)\rfloor\)
and \(B=4Lb\). These choices give \(b\ge2,\ b\ge e^{S_0/d_0}/(8L)\).
Encode \(\{0,\ldots,b-1\}^{d_0}\) as base-\(B\) integers. All encoded
values are \(<B^{d_0}\le n/(4L)\). One level of the squared Euclidean norm
has at least \(b^{d_0}/(2d_0 b^2)\) vectors. Its log cardinality is at least
\[
 S_0-d_0\log(8L)-2S_0/d_0-\log(2d_0)
 \ \ge\ S_0-C_1\sqrt{S_0 v}.
\]
Here \(d_0\) lies between \(\tfrac12\sqrt{S_0/v}\) and \(\sqrt{S_0/v}\).
Absorbing \(\log4\) and using \(v\le3\log L\) gives the claimed size.
If the asserted integer relation held, each digit of the signed
difference would have absolute value at most \(2L(b-1)<B\).
Reducing successively modulo \(B\) forces every such digit to be zero.
This would give three distinct collinear vectors on the same Euclidean
sphere, which a line meets in at most two points.

For three hits at times \(a<b'<c\) in a length-\(L\) segment with step
of order \(>4L\), the representatives \(s_1,s_2,s_3\) are distinct.
The translates cancel, so
\((c-a)(s_2-s_1)-(b'-a)(s_3-s_1)\equiv0\bmod n\).
Its absolute value is \(<n\), and its vanishing over \(\mathbf Z\)
contradicts the preceding property. For an orbit of order \(m'\le4L\),
any two hits have \(n\mid m'(s_2-s_1)\) with absolute value \(<n\).
They therefore coincide.
\end{proof}

In the application \(n_0=|A_0|=n/o\), \(o\) is fixed and
\(L_w=10\lceil405\log x\rceil+2\lceil30/\varepsilon\rceil+22\).
Use \(S\) in the cyclic group \(A_0\) with \(n_0,L_w\) in place of
\(n,L\). For fixed \(\varepsilon\) its lower bound divided by \(n\)
is at least \(x^{-\varepsilon}\) eventually: the logarithm of its
reciprocal is bounded by
\(\log(oL_w)+C\sqrt{\log n_0\log L_w}=o(\log x)\).
Thus the set giving the lower density also supplies the orbit restriction
at every translate.

\begin{lemma}[One choice of the top set]\label{host:shares}
In a finite group \(H\) of order \(n\asymp x^{2+8\varepsilon}\) with
subgroup \(A_0\) of fixed index \(o\), put \(\delta=x^{-\varepsilon/8}\).
Include points independently with probability \(\delta\) on the identity
coset and \(3\delta/4\) on the others. Let \(D(T)\ge0\) be any weighted
cost with \(\mathbf E D(T)\le B_0\). For sufficiently large \(x\) there is
one \(T\) with \(D(T)\le4 B_0,\ |T|\ge n\delta/2\), and
\[
  \big||T\cap C|-\lambda_C|T|\big|\le |T|x^{-1+\varepsilon}
  \quad\text{at every coset},\qquad
  \lambda_C=\frac{3+\mathbf1_{C=A_0}}{3o+1}.
\]
\end{lemma}
\begin{proof}
Let \(\mu=\mathbf E|T|\ge3n\delta/4\). Each coset count has mean
\(\lambda_C\mu\) and variance at most \(\mu\), by independence. For
\(a=\mu x^{-1+\varepsilon}/(4(o+1))\), Chebyshev and a union bound
give failure probability at most \(o\mu/a^2=O_o(x^{2-2\varepsilon}/\mu)
=o(1)\). On success, \(\big||T|-\mu\big|\le oa\), hence
\(|T|\ge n\delta/2\) and \((o+1)a\le |T|x^{-1+\varepsilon}\).
Apply the triangle inequality to each coset. Markov's inequality gives
\(\Pr(D>4B_0)\le1/4\) for \(B_0>0\); for \(B_0=0\) the cost is zero
almost surely. The two good events intersect.
\end{proof}

For example, if the cost counts chosen points \(t\) for which at least
\(k+1\) points of some specified window \(W\) of size at most \(L\) are
chosen, its expectation is bounded by a union over the windows and their
subsets. For \(t\in W\), the event requires \(t\) and at least \(k\)
other points, costing at most \(\delta(L\delta)^k\); for \(t\notin W\),
the bound is \(\delta(L\delta)^{k+1}\). This supplies a valid input
\(B_0\) whenever the windows and their weights are specified.
Lemma~\ref{host:shares} does not give a simultaneous automorphic realization
of the resulting translates. That additional assertion identifies the
centre \(a_t\) and the weights \(w_{u,j,\rho}\) for all the local periods
with the \emph{same} values in the Fourier discrepancy. The quantitative
form still needed is given in Section~\ref{scope:account}.

\subsection{Bounds for all characters at the thin coordinates}

\begin{proposition}[Uniform thin-coordinate estimate]\label{host:thin}
Fix \(M,\varepsilon\) as in Theorem~\ref{pr:pair}. Suppose
\(Y\le x\), \(x^{2+8\varepsilon}\le |H_e|\le2x^{2+8\varepsilon}\),
and \(q,d\) are the pair selected at \(Y\). For all sufficiently large
\(x\) (with \(Y\) above the pair threshold), the primitive Dirichlet inducer
\(\chi\) of each character of \(\operatorname{Gal}(E/\mathbf Q)\), for the field \(E\) of Proposition~\ref{host:art}, satisfies
\[
        |G_\chi(y)|\le y e^{-x^{\varepsilon/8}}
        \qquad \bigl(y\ge \exp(x^{\varepsilon/3})/3\bigr).
\]
\end{proposition}
\begin{proof}
Write \(N\) for the primitive conductor. From \(N\mid N_eqd\),
\(2^e\le|H_e|\), and the ranges of the pair, we obtain
\(H=\log(2N)\ll_{M,\varepsilon}x^{\varepsilon/5}(\log x)^2\).
There is a fixed \(\nu>0\) such that any exceptional real zero of \(\chi\),
if present, is at most \(1-\nu\). A primitive real character
on the cyclotomic modulus \(N_eqd\) factors as a product of local
quadratic or trivial characters. At an odd prime the principal units
are a group of odd order and are killed by a quadratic character.
At 2 every odd square is \(1\bmod8\), and every unit \(1\bmod8\) is
a square modulo a power of 2 (lift an odd square one bit at a time).
The part away from \(q,d\) therefore has conductor dividing \(8M\).
If a \(q\)- or \(d\)-factor is present, Theorem~\ref{pr:pair} supplies
the gap \(1/400\) at real height zero. The remaining real primitive
characters form a fixed finite set. Classical nonvanishing at one and
analyticity give them a common positive gap near one. At conductor one
use analyticity and nonvanishing of \((s-1)\zeta(s)\) at one.

Let \(D_\chi(t)=\sum_{n\le t}\Lambda(n)\chi(n)
-\mathbf1_{\chi=1}t\).
For \(t\ge(2N)^{12}\) sufficiently large, Corollary~1.4 of \cite{tz}
applied at modulus \(N\ge2\) gives
\[
 |D_\chi(t)|/t\ll
 e^{-c\log t/H}+e^{-c\sqrt{\log t}}+t^{-\nu}+t^{-1/4}.
\]
Indeed, multiply the formula for each unit class by \(\chi\) and sum.
Its primary terms cancel for a nonprincipal character. Its exceptional
term cancels unless \(\chi\) is the corresponding real character, in
which case the real gap bounds it. The absolute errors add to
\(O(t\Delta)\) because \(H_b(t)\le 3t/\varphi(N)\) in the notation
of Lemma~\ref{pr:signs}. Prime powers of exponent at least two contribute
\(O(\sqrt t\log^2t)\), by summing the trivial bound for their prime bases.
The second exponential weakens the corresponding bound in Corollary~1.4.
For \(N=1\), use the same formula for the odd class modulo two, add the
prime \(2\) and then the higher powers.

Summation of the ramp gives
\(P_\chi(y)=y\int_1^y(\sum_{n\le t}\Lambda(n)\chi(n))t^{-2}\,dt\).
In fact, for both kinds of character,
\[
            G_\chi(y)=-y\int_y^\infty D_\chi(t)t^{-2}\,dt .
\]
To identify the constant, use
\(-L'/L(s,\chi)=s\int_1^\infty
(\sum_{n\le t}\Lambda(n)\chi(n))t^{-1-s}\,dt\) for real \(s>1\).
After subtraction, the displayed error is integrable against \(dt/t\),
so the limit \(s\downarrow1\) gives the nonprincipal constant. In the
principal case subtract \(s/(s-1)\) before taking the limit; the
result is \(-1-\gamma\), as required by \eqref{eq:centred}.

Put \(A=\log y\ge x^{\varepsilon/3}-\log3\). Integrating the first two
errors from \(A\) to infinity in the variable \(\log t\) gives
\(O(H e^{-c A/H}+(1+\sqrt A)e^{-c\sqrt A})\);
the power errors give \(O_\nu(e^{-\nu A}+e^{-A/4})\).
All are \(o(e^{-x^{\varepsilon/8}})\) uniformly up to \(y=\infty\).
Indeed \(A/H\gg x^{2\varepsilon/15}/(\log x)^2\) and
\(\sqrt A\gg x^{\varepsilon/6}\), with both exponents strictly larger
than \(\varepsilon/8\). In particular \(y\) is eventually above
\((2N)^{12}\).
\end{proof}

\begin{lemma}[Variation of the ramp]\label{host:variation}
For \(a,w,b>0\), extending the same ramp formula to \(0<y<1\) where
necessary, one has
\[
 |wG_\chi(a/w)-bG_\chi(a/b)|\le 3a\,|w-b|/\min(w,b).
\]
For \(y,y'\ge1\) the change in \(G_1\) is \(O(|y-y'|)\); that in
\(G_\chi\) for \(\chi\ne1\) is at most
\(|y-y'|(\log\max(y,y')+C+|L'(1,\chi)/L(1,\chi)|)\).
\end{lemma}
\begin{proof}
Write \(\Psi(v)=\sum_{n\le v}\Lambda(n)\). One has
\(\Psi(v)\le 4(\log2)v<3v\) also for small \(v>0\).
Indeed each prime power between \(2^{j-1}\) (excluded) and \(2^j\)
contributes its prime to \(\binom{2^j}{2^{j-1}}\) by the factorial
valuation formula. Bound the binomial coefficient by \(2^{2^j}\), sum
over \(j\), and round \(v\ge1\) up to a power of two. Below one
\(\Psi(v)=0\). Off the finitely many breakpoints on a compact interval
of \(u>0\), the derivative of
\(uP_\chi(a/u)=\sum_n\Lambda(n)\chi(n)(a/n-u)_+\)
has absolute value at most \(\Psi(a/u)\).
The nonprincipal compensation is constant after multiplication by \(u\).
In the principal case the total derivative is \(a/u-\Psi(a/u)\),
also of magnitude at most \(3a/u\). Integrate between \(w\) and \(b\).
For variation in \(y\), differentiate the ramp itself and use
\(\sum_{n\le y}\Lambda(n)/n=\log y+O(1)\), proved in
Theorem~\ref{pr:pair}. For the principal term subtraction of \(\log y\)
leaves a bounded derivative. The functions are continuous at breakpoints,
so integration gives the asserted bounds there as well.
\end{proof}

\begin{proposition}[The analytic CM scalar on the selected hosts]\label{host:period}
Fix a primitive character \(\psi\)
and a positive exponent \(\varepsilon\) as in
Theorem~\ref{pr:pair}. At a selected scale \(x\) let \(E\) be the
field in Proposition~\ref{host:art}, let \(d\) be the second
auxiliary prime and put \(L=x^\varepsilon\).
In the normalization of the determinant comparison in Section~9 the
real scalar is
\[
 A_{J_d}=\gamma+\log(2\pi)-\log\operatorname{rd}(E)
                         -\frac{L'(1,J_d)}{L(1,J_d)} .
\]
Uniformly in all allowed primes and cyclotomic levels at \(x\),
\[
 |A_{J_d}|+\log d+\log\operatorname{rd}(E)
          \ll_\psi (1+x^{\varepsilon/5})\log^2(2x)=o(L).
\]
Here and in the little-oh estimate the threshold may depend on
\(\psi,\varepsilon\). In particular the analytic size condition
\(|A_{J_d}|\le3L\) holds throughout every sufficiently late selected
window.
\end{proposition}
\begin{proof}
Write \(M\) for the primitive conductor of \(\psi\).
At the page used in the construction \(2^e\le2x^{2+8\varepsilon}\).
The field \(E\) is contained in the cyclotomic field of conductor
dividing \(C=\operatorname{lcm}(2^{e+2},M)qd\).
The discriminant tower gives \(\operatorname{rd}(E)\le
\operatorname{rd}(\mathbf Q(\zeta_C))\le C\). For the second
inequality, factor \(X^C-1=\Phi_C(X)Q(X)\), differentiate at
\(\zeta_C\), and take norms. Then
\(|N\Phi_C'(\zeta_C)|\le C^{\varphi(C)}\), since \(Q(\zeta_C)\)
is a nonzero algebraic integer and \(\mathbf Z[\zeta_C]\) is the
integer ring. On the initial interval starting at \(Y_*\le x\),
the prime bounds are
\(\log q\le Y_*^{\varepsilon/5}\) and
\(\log d\le Y_*^{\varepsilon/5}(\log Y_*)^2\).
They give the stated bound for \(\log C\), including the dyadic
level and the fixed conductor. The quadratic character of the
imaginary field of discriminant \(-d\) is \(J_d\), since \(d\)
is prime and \(d\equiv7\bmod16\). The formula for \(A_{J_d}\)
is the definition of the analytic scalar, and all its terms are
real. Apply Theorem~\ref{pr:pair} to the trivial base character
and its \(d\)-twist to obtain
\(|L'(1,J_d)/L(1,J_d)|\le4\log x\). Finally
\(x^{-4\varepsilon/5}\log^2(2x)\to0\).
\end{proof}
Consequently, wherever the determinant-degree comparison of
Section~9 identifies the corrected constant as
\[
 c_1=-\tfrac12(p-q')A_{J_d}
                +q'\big(\log\operatorname{rd}(E)-\log d\big),
\]
one has \(|c_1|=o(rL)\) uniformly even for varying signature
\((p,q')\), \(r=p+q'\ge1\). This removes a numerical period bound
from the list of hypotheses to be checked on the actual hosts.
Identifying an integral determinant degree with that scalar is a
geometric assertion.
\src{src/Development/AbelianVarietyTheory/Definitions/CMLattices/ExistsIsThetaDegOfIsLocalModel.lean}{23}{normalization of the analytic CM scalar} and
\src{src/Development/AutomorphicArithmetic/CharacterBoxes/ExistsPageN.lean}{39}{dyadic level at the selected scale}.
